% SPDX-License-Identifier: CC-BY-4.0
% arXiv preprint source (LaTeX). Target category: cs.CR (Cryptography and
% Security); cross-list cs.CY (Computers and Society). Content mirrors
% ep-multi-handshake-preprint.md in this directory.
% Options for packages loaded elsewhere
\PassOptionsToPackage{unicode}{hyperref}
\PassOptionsToPackage{hyphens}{url}
\documentclass[
  11pt,
]{article}
\usepackage{xcolor}
\usepackage[margin=1in]{geometry}
\usepackage{amsmath,amssymb}
\setcounter{secnumdepth}{-\maxdimen} % remove section numbering
\usepackage{iftex}
\ifPDFTeX
  \usepackage[T1]{fontenc}
  \usepackage[utf8]{inputenc}
  \usepackage{textcomp} % provide euro and other symbols
\else % if luatex or xetex
  \usepackage{unicode-math} % this also loads fontspec
  \defaultfontfeatures{Scale=MatchLowercase}
  \defaultfontfeatures[\rmfamily]{Ligatures=TeX,Scale=1}
\fi
\usepackage{lmodern}
% Use upquote if available, for straight quotes in verbatim environments
\IfFileExists{upquote.sty}{\usepackage{upquote}}{}
\IfFileExists{microtype.sty}{% use microtype if available
  \usepackage[]{microtype}
  \UseMicrotypeSet[protrusion]{basicmath} % disable protrusion for tt fonts
}{}
\makeatletter
\@ifundefined{KOMAClassName}{% if non-KOMA class
  \IfFileExists{parskip.sty}{%
    \usepackage{parskip}
  }{% else
    \setlength{\parindent}{0pt}
    \setlength{\parskip}{6pt plus 2pt minus 1pt}}
}{% if KOMA class
  \KOMAoptions{parskip=half}}
\makeatother
\usepackage{longtable,booktabs,array}
\usepackage{calc} % for calculating minipage widths
% Correct order of tables after \paragraph or \subparagraph
\usepackage{etoolbox}
\makeatletter
\patchcmd\longtable{\par}{\if@noskipsec\mbox{}\fi\par}{}{}
\makeatother
% Allow footnotes in longtable head/foot
\IfFileExists{footnotehyper.sty}{\usepackage{footnotehyper}}{\usepackage{footnote}}
\makesavenoteenv{longtable}
\setlength{\emergencystretch}{3em} % prevent overfull lines
\providecommand{\tightlist}{%
  \setlength{\itemsep}{0pt}\setlength{\parskip}{0pt}}
\usepackage{bookmark}
\IfFileExists{xurl.sty}{\usepackage{xurl}}{} % add URL line breaks if available
\urlstyle{same}
\hypersetup{
  pdftitle={Handshakes, Not Signatures: Composing Fresh, Device-Bound Approval Ceremonies into Accountable M-of-N Authorization},
  pdfauthor={Iman Schrock (EMILIA Protocol, Inc.; team@emiliaprotocol.ai)},
  hidelinks,
  pdfcreator={LaTeX}}

\title{Handshakes, Not Signatures: Composing Fresh, Device-Bound Approval
Ceremonies into Accountable M-of-N Authorization}
\author{Iman Schrock (EMILIA Protocol, Inc.; team@emiliaprotocol.ai)}
\date{5 July 2026}

\begin{document}
\maketitle

\emph{Preprint, 5 July 2026.} Companion to the EP-QUORUM preprint
(Zenodo, DOI
\href{https://doi.org/10.5281/zenodo.20780638}{10.5281/zenodo.20780638})
and to IETF Internet-Draft
\href{https://datatracker.ietf.org/doc/draft-schrock-ep-quorum/}{draft-schrock-ep-quorum}.
That paper specifies the quorum predicate and its conformance vectors.
This paper argues the design decision underneath it: what the unit of
multi-party human authorization should be, and why.

\subsection{Abstract}\label{abstract}

When a system requires M of N humans to authorize an irreversible
action, it must choose a unit of composition. The cryptographic
literature offers the signature: threshold and aggregate schemes (FROST,
MuSig2, threshold ECDSA) combine signing shares into one compact
signature, and collect-then-verify constructions (CMS multi-signer,
script-level multisig) gather independent signatures over the same
bytes. We argue that for \emph{human} authorization of high-consequence
machine actions, the signature is the wrong unit, and that schemes built
on it commit a category error: they optimize for compactness and key
custody where the problem demands evidence and accountability. We
propose the \emph{handshake} as the unit of composition: a fresh,
intent-bound, user-verified challenge-response ceremony with an explicit
lifecycle and a recorded terminal outcome, including refusal. A
multi-party authorization is then an ordered trail of handshakes, each
admitted against the governing policy \emph{before} it is recorded,
optionally chained so that each approver signs over a commitment to its
predecessor, and consumed exactly once as a composed decision. We
identify five properties that handshake composition preserves and
signature aggregation destroys (per-member freshness, per-member intent
binding, individual accountability including denials, order evidence,
and machine-checkable separation of duties), and for each we name the
attack that exists in its absence. The model's safety core is
machine-checked: a TLC run over the single-handshake state machine,
exhaustive at its bounded-exploration configuration (413,137 states, 26
invariants, no counterexample), a bounded two-handshake exploration, and
Alloy checks of the composition
predicate, with cryptographic operations treated as axioms at that
level; and, going beyond the axiom for the base receipt, a first symbolic
Dolev-Yao model in Tamarin verifies core authenticity and shows that
removing the one-time-consumption check makes the prover construct the
replay attack itself. The process also emits ceremony telemetry that lets an evidence
layer downgrade rubber-stamped approvals. We state plainly what no
composition of ceremonies can solve: collusion of the required number of
humans, coercion, person-level Sybil enrollment, and the semantic
correctness of the approved action.

\subsection{1. Introduction}\label{introduction}

Multi-party authorization is old. Banks have required two signatures on
large instruments for centuries; weapons handling and treasury movement
formalized the two-person rule; modern change management requires a
reviewer who is not the author. When these controls are ported to
software, and especially to the authorization of actions proposed by AI
agents, an engineering question appears that the paper forms never had
to answer: \emph{what, exactly, is the thing we collect M of?}

The available cryptographic answer is the signature. A rich literature
composes signatures: threshold schemes derive one signature from M-of-N
signing shares; multisignature schemes aggregate N contributions into
one compact artifact; plainer constructions simply collect N independent
signatures over the same message and verify each. All of these answer
the question ``did enough keys participate?'' That is the right question
for key custody, where the signers are machines or the goal is to make a
group of signers look like one. It is the wrong question for human
authorization of an irreversible action, where the artifact that
survives the action must support a different interrogation: \emph{which
accountable people said yes, to exactly what, when, in what order,
having seen what, and who said no?}

This paper makes that argument concretely, against a shipped system. The
EMILIA Protocol (EP) binds a named human approver to one exact action
through a challenge-response ceremony gated by user verification on the
approver's own device; the quorum profile (EP-QUORUM) composes several
such ceremonies into an M-of-N or ordered authorization {[}1, 2, 3{]}.
The companion paper {[}3{]} specifies the composition predicate and its
adversarial conformance vectors. Here we defend the design decision
beneath the predicate:

\begin{enumerate}
\def\labelenumi{\arabic{enumi}.}
\tightlist
\item
  \textbf{The unit of multi-party human authorization should be the
  handshake}, defined as a fresh, intent-bound, user-verified
  challenge-response ceremony with an explicit lifecycle whose terminal
  outcomes include refusal (Section 3).
\item
  \textbf{Composition of handshakes preserves five properties that
  aggregation of signatures destroys}, each of which corresponds to a
  concrete attack when absent (Section 4).
\item
  \textbf{The safety core of the composed process is machine-checkable
  at the state-machine level}, and we report exactly what has been
  checked, at what scope, and what has not (Section 5).
\item
  \textbf{The ceremony yields telemetry that makes rubber-stamping
  visible}, which signature-only constructions structurally cannot
  (Section 6).
\item
  \textbf{Composition has limits no process design removes}, and we
  state them (Section 7).
\end{enumerate}

A terminology note for readers with the EP repository open. EP's
codebase uses ``handshake'' for the verified counterparty agreement
object that every approval binds into, and ``signoff'' for the
per-approver challenge-response ceremony (challenge, attestation,
terminal outcome). In this paper ``handshake'' means the ceremony,
because the word names precisely what we argue the unit should be: a
live exchange with a specific person, not a static artifact from one.
Where we cite code or invariants we use the repository's names.

\subsection{2. The Category Error}\label{the-category-error}

Consider the design space of ``M humans must approve'' as it is usually
built.

\subsubsection{2.1 Aggregating schemes produce one
signature}\label{aggregating-schemes}

Threshold signature schemes (threshold ECDSA {[}6{]}, FROST {[}4, 5{]})
and multisignature schemes (MuSig2 {[}7{]}) share an objective: N
parties hold key material, some quorum of them cooperates, and the
output is \emph{one} signature, verifiable against \emph{one} public
key, indistinguishable from a single-signer signature. This is a
feature, and a hard-won one: compact verification, constant-size
artifacts, no on-chain fee growth, no revelation of the signer set or
the policy.

Every one of those features is an anti-feature for human authorization.

\begin{itemize}
\tightlist
\item
  \textbf{Attribution is destroyed by construction.} The verifier learns
  that some authorized quorum signed; it cannot learn which members.
  Accountable-subgroup multisignatures {[}8{]} recover attribution, but
  as an opt-in property grafted onto a primitive whose default erases
  it.
\item
  \textbf{Freshness does not survive into the artifact.} These protocols
  use per-session nonces internally (nonce reuse is catastrophic for
  them), but that freshness is machinery between co-signers. The final
  signature carries no per-member evidence of \emph{when} each human
  participated or that each participation was contemporaneous with this
  decision.
\item
  \textbf{There is no ceremony.} The protocol steps are executed by
  software holding shares. Nothing in the output attests that a human
  saw a rendering of the action and made a verification gesture. A share
  held on a server signs exactly as convincingly as a share held behind
  a fingerprint.
\item
  \textbf{Refusal does not exist.} A member who declines simply does not
  participate. The artifact of a 2-of-3 signing where the third member
  objected strenuously is byte-identical to one where the third member
  was on vacation.
\item
  \textbf{Order does not exist.} One signature has no internal sequence.
\end{itemize}

\subsubsection{2.2 Collect-then-verify keeps attribution, and little
else}\label{collect-then-verify}

The plainer construction gathers N independent signatures over the same
bytes: CMS with multiple SignerInfo structures {[}9{]}, script-level
M-of-N multisig as deployed in Bitcoin, and the multi-signature approval
flows of on-chain wallet contracts. Attribution survives: each signature
verifies against a named key. This is materially better, and it is the
branch of the design space EP-QUORUM descends from. But as usually built
it still lacks:

\begin{itemize}
\tightlist
\item
  \textbf{Per-signer freshness.} Each member signs the same static
  message. The scheme cannot distinguish a signature produced this
  morning in review of this decision from one produced last year over
  the same bytes, harvested and replayed, or pre-signed in bulk.
  (Transaction digests give action binding, so cross-action replay
  fails; within the action, nothing binds the time or occasion of any
  member's consent.)
\item
  \textbf{An intent ceremony.} A key signed; whether a person reviewed
  anything is outside the model. Signing-time attributes, where they
  exist (CMS), are signer-asserted rather than challenge-anchored.
\item
  \textbf{Ordering evidence.} Signatures are an unordered bag, or
  ordered by the assembler. Where an order is checked (script
  evaluation), it is the assembly order of the presented artifact, not
  evidence of the order in which humans decided.
\item
  \textbf{Separation-of-duties enforcement.} Distinctness of
  \emph{keys} is checkable; distinctness of \emph{humans} is not
  expressible, nor is ``the proposer of this action must not be among
  its approvers.''
\item
  \textbf{A record of refusal.} As above: absence, not evidence.
\end{itemize}

\subsubsection{2.3 The trade human authorization actually
needs}\label{the-trade}

The common root is that both branches treat the signature as the atom
and therefore inherit the signature's semantics: a signature proves that
a key participated in a computation over bytes. Everything the auditor
of a high-consequence action needs beyond that (who, when, in view of
what, in what order, over whose objection) must either be bolted on or
is unrecoverable.

Human authorization needs the opposite trade from the one aggregation
makes. Compactness is nearly worthless: an authorization for an
irreversible wire executes once, and the artifact is read by auditors,
not merged into blocks. Evidence is nearly everything. The composed
artifact should therefore be \emph{the concatenation of everything each
member did}, with the composition rules enforced over it. Its size is
linear in M, and that is not a cost to be engineered away; the size
\emph{is} the evidence.

We now make ``everything each member did'' precise.

\subsection{3. The Multi-Handshake
Model}\label{the-multi-handshake-model}

\subsubsection{3.1 The handshake}\label{the-handshake}

Fix a canonical encoding and a hash function (EP uses JCS-style
canonicalization {[}10{]} and SHA-256). Let $h_A$ be the digest of the
exact action bytes to be authorized and $h_P$ the digest of the
governing policy, including the quorum policy of Section 3.2.

A \textbf{handshake} is a tuple
\[
H = (a,\; k_a,\; c,\; \sigma,\; s)
\]
where:

\begin{itemize}
\tightlist
\item
  $a$ is an enrolled approver identifier, and $k_a$ the device-bound
  public key enrolled for $a$ out of band (server-side enrollment; never
  taken from the artifact under verification);
\item
  $c$ is the \textbf{challenge context}:
  $c = (n,\, h_A,\, h_P,\, a,\, t_{\mathrm{iss}},\, t_{\mathrm{exp}}\, [,\, h_{\mathrm{prev}}])$,
  with $n$ a single-use nonce minted for this ceremony,
  $t_{\mathrm{iss}}$/$t_{\mathrm{exp}}$ the issuance and expiry
  instants, and $h_{\mathrm{prev}}$ an optional commitment to a
  predecessor context (Section 3.3);
\item
  $\sigma$ is a signature by $k_a$ over $c$ produced through a
  user-verified device ceremony: in EP, a WebAuthn assertion {[}11{]}
  whose challenge equals the hash of the canonical context and whose
  authenticator data asserts user verification, so the signature attests
  key possession \emph{and} a contemporaneous human-verification event
  on the approver's own device;
\item
  $s$ is the lifecycle state. The lifecycle is a forward-only order
\end{itemize}

\begin{verbatim}
challenge_issued -> challenge_viewed -> approved -> consumed
\end{verbatim}

with alternate terminal outcomes \texttt{denied}, \texttt{expired}, and
\texttt{revoked} reachable from the non-terminal states. The terminal
set is \texttt{\{denied,\ consumed,\ expired,\ revoked\}};
\texttt{approved} is deliberately \emph{not} terminal, because an
approval that has not been consumed authorizes nothing yet and can still
expire or be revoked. These states and their ordering are the code
constants \texttt{SIGNOFF\_STATUS\_ORDER} and
\texttt{SIGNOFF\_TERMINAL\_STATES} in
\texttt{lib/signoff/invariants.js}, and backward transitions are
excluded both in code and in the model (Section 5).

Three aspects of this definition carry the paper's argument. First, the
challenge is \emph{minted}, not derived: freshness is a property of the
ceremony, anchored by $n$ and bounded by $t_{\mathrm{exp}}$, rather
than an inference from signer-asserted times. Second, the signature is
\emph{gated}, not merely computed: the user-verification requirement
makes ``a human performed a deliberate gesture over a rendering of this
exact action'' part of what verifies. Third, the outcome is
\emph{total}: a handshake ends in exactly one of approval consumed,
denial, expiry, or revocation, so ``nothing happened'' is not a
representable result of asking a human for authorization.

\subsubsection{3.2 Trails and incremental
admission}\label{trails-and-incremental-admission}

A \textbf{quorum policy} $P$ declares a mode (\texttt{threshold} or
\texttt{ordered}), a required count, a roster of eligible
\texttt{(role,\ approver)} slots, a distinct-humans rule (default true),
an approval window, and optionally the strong ordering flag of Section
3.3. A \textbf{trail} for action $h_A$ under $P$ is an ordered sequence
of handshakes
\[
T = \langle H_1, H_2, \ldots, H_m \rangle
\]

The defining process rule is \textbf{incremental admission}: a candidate
handshake is checked against $(P, h_A, T)$ \emph{before} it is
recorded, and a candidate that fails any check never enters the trail.
The admission predicate $\mathrm{Admit}(P, h_A, T, H)$ requires, in
order: a well-formed policy with a non-empty roster; action binding
($c.h_A$ equals the trail's $h_A$); roster admission of the
\texttt{(role,\ a)} pair; human distinctness against the existing trail;
in ordered mode, that $H$ fills exactly the next unfilled roster slot;
window containment and, in ordered mode, strictly increasing signature
times; and verification of $\sigma$ itself. Each failure is a distinct
machine-readable refusal: \texttt{no\_policy},
\texttt{no\_eligible\_approvers}, \texttt{action\_mismatch},
\texttt{ineligible\_role}, \texttt{duplicate\_human},
\texttt{out\_of\_order}, \texttt{window\_exceeded},
\texttt{non\_increasing\_time}, \texttt{invalid\_signature}. This is the
\texttt{canAccept} function of \texttt{lib/signoff/quorum-session.js},
normatively specified in Section 6 of the quorum draft {[}2{]}.

Admission-before-recording is a process property, not an optimization,
and it is the first place the handshake model diverges operationally
from collect-then-verify. In a collect-then-verify system the trail is
whatever bag of signatures accumulated, and verification is a post-hoc
filter; malformed, duplicate, or off-roster contributions exist in the
record and must be argued about later. Under incremental admission
\textbf{a malformed trail never exists}: every recorded prefix of $T$
satisfies the policy's structural constraints by construction, the
refusal of a non-conforming signer is itself a logged, attributable
event, and at decision time there is no ambiguity about which members
``count.''

Admission is nonetheless \emph{not} trusted. The executing system
re-evaluates the full composition predicate $\mathrm{Sat}(P, h_A, T)$
over the assembled trail before acting, because it does not trust the
orchestrator to have admitted honestly. $\mathrm{Sat}$ is the
fail-closed quorum gate specified in {[}2{]} Section 5 and analyzed in
{[}3{]}: all member signatures verify with challenge equal to context
hash and user verification asserted; every member binds $h_A$;
approvers pairwise distinct and distinct from the initiator; device keys
pairwise distinct; every \texttt{(role,\ a)} on the roster; the
threshold met; the order and, when required, the ordering chain intact;
all signatures within the window. Any failure, malformation, or
unrecognized condition yields ``not satisfied.''

One discipline of the reference verifier deserves restating here because
it is about process rather than predicate: \textbf{verified is not
authorized-by-the-org}. $\mathrm{Sat}$ is a pure function of the policy
and keys it is handed. A creator-chosen policy verified against
creator-declared keys proves internal consistency only. The verifying
party must source $P$ from material it pins (in EP, an org-pinned
template keyed by organization and action type) and each $k_a$ from
server-side enrollment, never from the artifact under verification. The
composition inherits the accountability of its inputs, not the other way
around.

\subsubsection{3.3 The ordering commitment}\label{the-ordering-commitment}

Ordered mode expresses escalation: a later authority signs after, and in
knowledge of, an earlier one. Timestamps alone cannot prove that,
because $t_{\mathrm{iss}}$ is stamped by whoever operates the clock.
The strong ordered mode (\texttt{ordered\_chain} in {[}2{]} Section 4)
therefore places the order \emph{inside the signed bytes}: each context
after the first carries $h_{\mathrm{prev}}$, the SHA-256 of the
canonical context of its predecessor, and the first carries none.
Because each approver's device signs its own context, the i-th approver
attests to having signed after the exact (i-1)-th approval. No party,
including the operator, can reorder, insert, or backdate a member
without invalidating a signature; the reference verifiers surface this
as \texttt{chain\_linked:\ false} (the specification names the refusal
\texttt{broken\_chain}), and a trail whose chain does not reconstruct is
not satisfied.

The chain also makes the trail \emph{linear} by construction: one root,
no branching, no cycles. These are exactly the properties checked in the
Alloy model (Section 5.2).

\subsubsection{3.4 Denial as a first-class
outcome}\label{denial-as-a-first-class-outcome}

A handshake whose human says no terminates in \texttt{denied}, an
attributed, event-logged, irreversible outcome: only the accountable
approver can deny their own challenge, the denial is recorded
event-first in the evidence log, and a denied ceremony can never later
be counted as approved (the machine-checked invariant
\texttt{DenyCannotBeApproved}, Section 5.1). The trail of a contested
authorization therefore contains the contest: an auditor reconstructing
a 2-of-3 approval sees the refusals alongside the approvals, with names
and times.

We are precise about current scope: in the implementation, denial is an
authenticated, logged decision by the accountable approver
(\texttt{lib/signoff/deny.js}); it does not yet carry a device-signed
assertion the way approval does. Extending denial to the same
device-signed evidence standard as approval is a natural hardening,
fully compatible with the model, and we note it as future work rather
than claiming it.

\subsubsection{3.5 One-time consumption of the composed
decision}\label{one-time-consumption}

A satisfied trail is not an authorization in perpetuity; it is
authorization for one execution. Per-member, an approved handshake is
consumed atomically at most once (attempting to consume a non-approved
or already-consumed attestation refuses; the invariants
\texttt{SignoffConsumeOnce} and \texttt{ConsumeOnceSafety} cover reuse).
At the composed level, the executing system's consume path refuses with
an explicit \texttt{quorum\_not\_satisfied} error until $\mathrm{Sat}$
holds over the assembled trail, and consuming terminalizes the
authorization. A partial trail confers no authority at any point in its
life: before the threshold it fails $\mathrm{Sat}$; after consumption
it is spent.

\subsubsection{3.6 Offline verifiability, honestly
scoped}\label{offline-verifiability}

Everything $\mathrm{Sat}$ needs (the policy, the members' contexts and
assertions, the pinned keys) is data. A third party can therefore
recompute the satisfied/not-satisfied judgment offline, without querying
the operator, and two correct verifiers cannot disagree on the same
inputs. We scope this claim carefully: offline verification establishes
that the composition was \emph{validly formed} relative to the pinned
policy and keys. It does not establish \emph{current} validity: whether
the authorization has since been consumed, or revoked, or the keys
withdrawn is state, and answering it requires a fresh view of that state
or an accompanying currency proof. A relying party that treats an
offline ``satisfied'' as ``currently valid'' has reintroduced the replay
problem one level up.

\subsection{4. What Composition Preserves and Aggregation
Destroys}\label{what-composition-preserves}

We now argue the five properties from Section 1. For each: the property
as the handshake model provides it, why signature-unit schemes lose it,
and the attack that exists when it is absent. Each property is pinned by
named adversarial conformance vectors from the EP-QUORUM-v1 suite
(\texttt{conformance/vectors/quorum.v1.json}, eleven vectors), which the
JavaScript, Python, and Go reference verifiers are required to agree on;
those are cross-language reference verifiers, one team's ports serving
as a consistency check, not clean-room independent implementations
(Section 9).

\subsubsection{4.1 Per-member freshness}\label{per-member-freshness}

\textbf{Provided.} Each member's signature exists only because a
challenge was minted for this ceremony: single-use nonce, bounded
validity, a challenge record that can be attested against only while
live, and a composed window constraint (\texttt{window\_sec}) bounding
the spread of the whole trail. A partial trail goes stale as a unit: the
window bounds how long a dormant partial quorum remains completable by
an attacker who compromises a remaining approver later.

\textbf{Lost.} Aggregation: internal nonces do not survive into the
artifact, so the verifier cannot check any member's freshness.
Collect-then-verify: signatures over static bytes are occasion-free;
nothing distinguishes fresh review from harvest or bulk pre-signing.

\textbf{Attack when absent: signature harvesting and replay across
ceremonies.} An operator (or an attacker holding the operator's
position) accumulates signatures produced on other occasions over
recurring action bytes, or extracts pre-signed approvals, and assembles
a ``quorum'' no member decided this time. In the handshake model the
replayed member fails at admission (a reused or expired challenge cannot
be attested; a member outside the window is refused) and again at
$\mathrm{Sat}$. Pinned by \texttt{reject\_expired\_window} and, for the
demo's stale-challenge refusal, the composition demo of Section 9.

\subsubsection{4.2 Per-member intent binding}\label{per-member-intent-binding}

\textbf{Provided.} What each member signs \emph{is} the intent object:
the challenge equals the hash of a canonical context that carries the
exact action digest and policy digest, and the signature is gated by
user verification on the member's own device. The artifact thus
witnesses, per member, ``a human performed a deliberate verification
gesture over a rendering bound to exactly these action bytes under
exactly this policy version.'' A signature collected under one policy
version cannot satisfy a requirement evaluated under another, because
$h_P$ is inside the signed bytes.

\textbf{Lost.} Aggregation: shares sign protocol messages; no human
gesture is part of what verifies. Collect-then-verify: the message is
bound but the reviewing human is a hope, not a property; a server-held
key is indistinguishable from an attended one.

\textbf{Attack when absent: consent laundering through blind signing.} A
key signs, and the system's owner asserts a human approved. Compromise
of the signing host silently converts ``two-person rule'' into
``two-key rule,'' which a single intruder with access to both hosts
satisfies alone. Pinned by \texttt{reject\_action\_mismatch} (a member
bound to different action bytes does not count) and
\texttt{reject\_one\_bad\_signature} (an assertion that does not verify
against the enrolled key fails the whole quorum).

\subsubsection{4.3 Individual accountability, including
refusal}\label{individual-accountability}

\textbf{Provided.} The trail names every member: identifier, role, key,
times, outcome. Approvals are device-signed; denials are attributed,
logged, and irreversibly terminal (Section 3.4). The composed artifact
is evidence \emph{for and against} each participant: an approver cannot
later deny approving, and an approver who refused cannot have their
refusal erased by re-asking, because the denied ceremony is terminal and
a fresh ceremony is a new, visible record.

\textbf{Lost.} Aggregation destroys attribution by design; accountable
variants {[}8{]} must re-add it. Collect-then-verify attributes
approvals but has no representation of refusal at all.

\textbf{Attack when absent: invisible refusals and approval shopping.}
An orchestrator polls approvers until it finds M willing; the artifact
shows M clean approvals and no trace of the N-M refusals, so governance
sees consensus where there was contest. Post-hoc deniability is the
mirror attack: with one aggregate signature, any individual member can
claim non-participation. The handshake model makes the first visible
(denials persist in the record) and the second impossible (each approval
verifies against one enrolled key).

\subsubsection{4.4 Order evidence}\label{order-evidence}

\textbf{Provided.} Ordered mode enforces roster order and strictly
increasing times at admission and at $\mathrm{Sat}$; strong ordered
mode proves order by the signatures themselves through the
$h_{\mathrm{prev}}$ chain (Section 3.3), with acyclicity and linearity
machine-checked (Section 5.2). ``The inspector general signed last, in
knowledge of both prior approvals'' is a verifiable statement, not an
operator log line.

\textbf{Lost.} One aggregate signature has no internal order. A
signature bag has whatever order the assembler asserts.

\textbf{Attack when absent: reordering and escalation forgery.} An
operator presents a junior's approval as having followed the senior's
(or the converse), fabricating an escalation chain that never occurred;
or backdates a later-obtained signature into an earlier slot. With the
chain, any such edit breaks a signature. Pinned by
\texttt{reject\_out\_of\_order} (roster order violated), the
non-increasing-time refusal at admission, and
\texttt{reject\_broken\_chain} (a final signoff committing to the wrong
predecessor hash fails even though every individual signature verifies).

\subsubsection{4.5 Machine-checkable separation of
duties}\label{machine-checkable-sod}

\textbf{Provided.} The composition rule, not organizational custom,
enforces: the initiator of the action is not a counted approver
(specified in {[}2{]} Section 5 check 5; modeled as
\texttt{SelfApprovalImpossible}; enforced at the write layer of the live
flow); approvers pairwise distinct (\texttt{duplicate\_human}); device
keys pairwise distinct, so one key enrolled under two identifiers cannot
fill two slots (\texttt{duplicate\_key}); and every member admitted
under an eligible \texttt{(role,\ approver)} slot, so a valid signature
by a real, enrolled, but off-roster approver is refused
(\texttt{wrong\_role}).

\textbf{Lost.} Threshold schemes cannot see humans at all: nothing
prevents one person from holding multiple shares, and the signature
cannot reveal it. Multisig can check key distinctness but cannot express
``distinct humans'' or ``not the proposer.''

\textbf{Attack when absent: self-approval and slot doubling.} The agent
(or its operator) that proposed the action supplies one or more of the
approvals; or one person, holding two keys or one key under two names,
satisfies a ``two-person'' rule alone. Pinned by
\texttt{reject\_duplicate\_human}, \texttt{reject\_duplicate\_key}, and
\texttt{reject\_wrong\_role}; threshold arithmetic itself by
\texttt{reject\_under\_threshold}; the two accept vectors
(\texttt{accept\_ordered\_3of3}, \texttt{accept\_threshold\_2of3}) guard
the other direction, a verifier too strict to accept valid compositions.

\subsubsection{4.6 The cost, stated plainly}\label{the-cost-stated-plainly}

The composed artifact is O(M) signatures plus contexts; verification is
M signature verifications plus set-level checks; the humans must
actually perform M ceremonies within a window. For machine-speed,
high-frequency, low-consequence signing, this is the wrong design, and
the aggregation literature serves those cases well. The claim of this
paper is scoped: for low-frequency, high-consequence, human-accountable
authorization, every one of those costs purchases exactly the property
the auditor needs, and the compact alternative discards evidence that
cannot be reconstructed afterward.

\subsection{5. Formal Grounding}\label{formal-grounding}

Machine-checked claims below name the exact identifiers as they appear
in the repository's \texttt{formal/} directory, with tool and scope. The
uniform honest scope: these are \textbf{state-machine level} results.
Cryptographic operations are axioms at this level; signature
unforgeability, WebAuthn internals, device binding, the approver
directory, log checkpoints, and wall-clock time are not modeled. A first
symbolic model in Tamarin (Section 5.3) goes further for the base
receipt: it verifies the core authentication lemma against a Dolev-Yao
adversary without axiomatizing the signature. Symbolic verification of
the full composed protocol (WebAuthn, directory, and log together)
remains future work, and we do not substitute the state-machine results
below for it.

\subsubsection{5.1 TLA+ / TLC: the handshake state machine,
exhaustively; two handshakes, bounded}\label{tla-tlc}

The model \texttt{formal/ep\_handshake.tla} (checked by TLC 2.19 under
\texttt{formal/ep\_handshake.cfg}) covers the full lifecycle of Section
3.1: challenge issuance against a verified counterparty agreement,
viewing, approval, denial, consumption, expiry, revocation, and the
delegation and policy machinery around them. The CI configuration
(\texttt{Handshakes\ =\ \{h1\}}, \texttt{Actors\ =\ \{a1,\ a2\}},
\texttt{Policies\ =\ \{p1\}}, \texttt{MaxPolicyVer\ =\ 2},
\texttt{Claims\ =\ \{c1\}}, exploration bounded to 10 events) is
\textbf{exhaustive at that configuration}: 413,137 states generated,
45,342 distinct states, depth 20, all 26 configured invariants hold with
no counterexample, re-run in CI on every push touching the TLA+ model or
its configuration.

The invariants that carry this paper's process claims, by their
identifiers in \texttt{ep\_handshake.cfg}:

\begin{longtable}[]{@{}
  >{\raggedright\arraybackslash}p{(\linewidth - 2\tabcolsep) * \real{0.45}}
  >{\raggedright\arraybackslash}p{(\linewidth - 2\tabcolsep) * \real{0.55}}@{}}
\toprule\noalign{}
\begin{minipage}[b]{\linewidth}\raggedright
Invariant
\end{minipage} & \begin{minipage}[b]{\linewidth}\raggedright
Process claim it grounds
\end{minipage} \\
\midrule\noalign{}
\endhead
\bottomrule\noalign{}
\endlastfoot
\texttt{SignoffRequiresVerifiedHandshake} & a ceremony exists only
against a verified counterparty agreement \\
\texttt{SignoffConsumeOnce} & a ceremony's approval is consumed at most
once (Section 3.5) \\
\texttt{SignoffBindingMatch} & a ceremony binds to the exact action
context (Section 4.2) \\
\texttt{SignoffTerminalIrreversible} & terminal outcomes are
irreversible (Section 3.1) \\
\texttt{DenyCannotBeApproved} & a denial can never become an approval
(Section 3.4) \\
\texttt{SignoffAuthorityMatch} & the signing actor is an authorized
approver (Section 4.5) \\
\texttt{ConsumeOnceSafety}, \texttt{ConsumeRequiresVerified} &
composed-decision consumption discipline \\
\texttt{WriteBypassSafety} & no approval-bearing write is reachable
except through the gate \\
\texttt{TerminalStateIrreversibility} & lifecycle is forward-only \\
\end{longtable}

Model checking earned its keep here in the ordinary way: TLC surfaced
four real specification bugs during verification (among them a
\texttt{Consume} action that could execute while a signoff was in
flight, and revocation/expiry that failed to cascade signoff
termination), which were fixed before the run above.

The \textbf{two-handshake configuration}
(\texttt{Handshakes\ =\ \{h1,\ h2\}}), which probes cross-ceremony
concurrency, was run under the same representation. It explored several
million distinct states with no counterexample before the search was
stopped: the state space explodes (tens of millions of states generated,
with the frontier still growing) and does not terminate at a CI-friendly
scale. Cross-handshake concurrency is therefore \textbf{bounded-tested,
not exhaustively proven}, and we claim exhaustive verification for no
configuration beyond one handshake. Exhaustively proving the composed
concurrent case is better suited to a symbolic prover than to
explicit-state enumeration.

\subsubsection{5.2 Alloy: the composition predicate and the relational
chain}\label{alloy}

\texttt{formal/ep\_quorum.als} models the composition predicate of
Section 3.2 directly (roster, initiator, members as signoffs with
holder-bound keys, an optional predecessor relation for the strong
chain) and declares six assertions, each checked at the scope recorded
in the file (\texttt{for\ 6}; \texttt{TwoPersonRuleHolds} at
\texttt{for\ 6\ but\ 4\ int}):

\begin{itemize}
\tightlist
\item
  \texttt{SelfApprovalImpossible}: a satisfied quorum never counts the
  initiator;
\item
  \texttt{NoHumanFillsTwoSlots}: pairwise-distinct approvers;
\item
  \texttt{NoKeyFillsTwoSlots}: one device key cannot fill two slots;
\item
  \texttt{TwoPersonRuleHolds}: a satisfied multi-member quorum contains
  at least two distinct humans;
\item
  \texttt{OrderedChainAcyclic}: the strong-chain predecessor relation
  has no cycles;
\item
  \texttt{OrderedChainLinear}: the chain is a single total order (one
  root, no branching), so the sequence cannot be permuted while
  remaining satisfied;
\end{itemize}

plus a non-vacuity \texttt{run} (\texttt{showStrongQuorum}) confirming a
genuine satisfied, chained, multi-member quorum exists in the model. The
CI workflow (\texttt{.github/workflows/alloy.yml}) checks every model in
\texttt{formal/*.als} with the SAT4J solver and fails the build on any
counterexample. Alloy results are bounded-exhaustive within the declared
scopes, not unbounded proofs.

\texttt{formal/ep\_relations.als} grounds the per-ceremony chain
relationally; the assertions relevant here, recorded in
\texttt{formal/PROOF\_STATUS.md} as verified at scope 6 with no
counterexample, are \texttt{SignoffBindingIntegrity} (A12, binding hash
consistent across handshake, challenge, attestation, consumption),
\texttt{SignoffConsumeOnce} (A13), \texttt{SignoffRequiresHandshake}
(A14), and \texttt{FullChainIntegrity} (A15).

\subsubsection{5.3 Tamarin: the core receipt lemma against a Dolev-Yao
adversary}\label{tamarin}

The models above treat the device signature as an uninterpreted verified
event. \texttt{formal/tamarin/ep\_receipt\_core.spthy} (checked by
tamarin-prover 1.10.0) does not: it runs a Dolev-Yao attacker who
controls the network and can ask the human to sign other actions, with an
explicit key-compromise rule, and models the user-verification event as a
linear fact a signature cannot exist without. The prover verifies
\texttt{executable\_honest\_receipt}, \texttt{core\_authenticity\_uv\_gated},
\texttt{no\_replay\_across\_actions}, and
\texttt{injective\_acceptance\_with\_consumption}, and \emph{falsifies}
\texttt{unchecked\_acceptance\_is\_injective}. Core authenticity holds: an
accepted receipt implies a user-verification-gated signature by the named
human over exactly that action, with no replay onto a different action.
The falsified lemma is by design: it asserts injective acceptance
\emph{without} the one-time-consumption check, and the prover constructs
the same-receipt replay trace itself (no forgery, no key reveal). That is
a mechanical demonstration that one-time consumption (Section 3.5) is
load-bearing, not defense-in-depth.

This model is scoped to the base single receipt. The approver directory,
log, checkpoints, WebAuthn attestation internals, and crucially the m-of-n
quorum and separation-of-duties composition of this paper are out of scope
here (they are carried by the state-machine models of 5.1 and 5.2).
Extending the symbolic model to the full composition is the next step,
and we claim no symbolic result beyond the base receipt.

\subsubsection{5.4 Enforced in code, not
model-checked}\label{enforced-in-code}

Two properties central to Sections 3.1 and 4.1 are deliberately
\emph{outside} the models and we do not count them as model-checked: the
user-verification / MFA gating of the attestation ceremony
(\texttt{lib/signoff/attest.js}, exercised by the WebAuthn signoff
conformance vectors) and challenge wall-clock expiry
(\texttt{lib/signoff/challenge.js}). Both are enforced by code guards
and covered by unit and conformance tests. The models treat time
abstractly and treat the device signature as an uninterpreted verified
event; that is the axiom boundary, stated rather than hidden.

\subsection{6. Ceremony Evidence and the Rubber-Stamp
Problem}\label{ceremony-evidence}

A known failure mode of multi-party approval is that it degenerates into
multi-party clicking. Adding signers can even reduce scrutiny, as each
assumes another is the careful one. Signature-unit schemes have nothing
to say here: a signature is a signature, produced in three seconds or
thirty minutes.

The handshake, because it is a lifecycle and not a point event, emits
telemetry as a side effect of operating: the challenge's
\texttt{issued\_at}, the \texttt{viewed\_at} transition, the
\texttt{approved\_at} instant, and the approver identity. EP's evidence
layer gives this a first-class node type, \texttt{ceremony\_evidence},
in the Action Evidence Graph (\texttt{lib/evidence/evidence-graph.js}),
and lets a relying-party policy declare a \textbf{review-latency floor}
(\texttt{ceremony\_min\_review\_sec}). The rule, shipped in that module,
is fail-closed in both directions:

\begin{itemize}
\tightlist
\item
  a verified ceremony node whose
  \texttt{approved\_at\ -\ viewed\_at} falls below the floor yields the
  machine-readable finding \texttt{rubber\_stamped\_ceremony} and
  downgrades the evaluation verdict to \texttt{conflicted}: the recorded
  human review contradicts a genuine one;
\item
  a floor that is set while the telemetry is absent or unusable yields
  \texttt{ceremony\_telemetry\_missing} and likewise downgrades to
  \texttt{conflicted}, never a silent pass;
\item
  no finding ever moves a verdict toward \texttt{admissible}, and an
  already-worse \texttt{unverifiable} verdict is never softened.
\end{itemize}

The claim is deliberately modest. A latency floor converts ``this
approver is at the floor every time'' from an anecdote into evidence
attached to the authorization itself, offline-checkable by any relying
party; it raises the cost of rubber-stamping and makes it attributable.
It does not prevent it. A colluding or coerced approver waits out the
floor. Detection, not prevention, is what process telemetry buys, and we
say so.

\subsection{7. What No Composition Solves}\label{what-no-composition-solves}

Composing ceremonies raises the cost of unilateral and unaccountable
action. Four problems survive any such composition, and a design that
claims otherwise is defective in a way this protocol treats as a
vulnerability class of its own: overclaiming discourages the
compensating controls that actually constrain these risks, and the base
specifications forbid it outright (implementations ``MUST NOT claim a
quorum is collusion-proof'' {[}2{]}).

\begin{enumerate}
\def\labelenumi{\arabic{enumi}.}
\tightlist
\item
  \textbf{Collusion of the required number of humans.} An M-person rule
  stops M-1 bad actors, not M. A fully colluding quorum produces a fully
  valid trail. What the trail changes is the aftermath: every colluder
  is named, signed, time-stamped, and action-bound, so the conspiracy is
  evidenced rather than deniable.
\item
  \textbf{Coercion.} A coerced approver's ceremony is cryptographically
  perfect. Ordering and windows narrow the operational space for
  coercing a full quorum simultaneously; they do not close it.
\item
  \textbf{Person-level Sybil enrollment.} If enrollment lets one human
  hold two approver identities on two devices, human-distinctness is
  satisfied formally but not in substance. The distinct-keys check
  closes the cheapest variant (one device key under two names,
  \texttt{reject\_duplicate\_key}); the general case is an
  enrollment-time control (identity proofing, directory authority), and
  no runtime predicate can supply it.
\item
  \textbf{Semantic correctness of the approved action.} M informed
  humans agreeing is the best available proxy for ``this action is
  right,'' and it is only a proxy. A quorum can competently, freshly,
  accountably approve a mistake. Rendering fidelity narrows the gap
  between what was approved and what was displayed; nothing in the
  composition addresses the gap between what was displayed and what was
  wise.
\end{enumerate}

\subsection{8. Related Work}\label{related-work}

Threshold and multisignature cryptography develops the signature-unit
branch: FROST {[}4{]}, standardized in RFC 9591 {[}5{]}; MuSig2 {[}7{]};
threshold ECDSA {[}6{]}. Accountable-subgroup multisignatures {[}8{]}
re-introduce signer attribution within aggregation and are the closest
cryptographic acknowledgment that accountability and aggregation pull
apart. CMS {[}9{]} defines the multi-signer collect-then-verify artifact
in wide deployment; script-level M-of-N multisig in Bitcoin and
multi-signature wallet contracts are its on-chain descendants. None of
these carries per-signer freshness ceremonies, intent gating, refusal
records, or order evidence, which is not a criticism of their design
goals; it is the reason a different unit is needed for human
authorization.

Two regulatory and standards artifacts already treat the ceremony, not
the key, as the unit of a human's authorization, and we read them as
convergent evidence. The DEA's electronic prescribing rule for
controlled substances (21 CFR Part 1311) requires the prescriber to
complete a two-factor authentication ceremony at the moment of signing
each prescription {[}12{]}: possession of the credential is explicitly
insufficient. OAuth 2.0 step-up authentication (RFC 9470) lets a relying
party demand a fresher, stronger authentication event for a specific
request {[}13{]}, acknowledging per-action freshness, though its result
is a session property rather than a portable, composable artifact.
WebAuthn {[}11{]} supplies the user-verified, device-bound,
challenge-response signature primitive the handshake is built from. The
EP base receipt {[}1{]} defines the single-ceremony artifact; the quorum
profile {[}2{]} and its companion paper {[}3{]} specify the composition
predicate this paper argues the foundations of.

\subsection{9. Artifact Availability}\label{artifact-availability}

The reference implementation is open source (Apache-2.0) in the EMILIA
Protocol repository (site: \url{https://www.emiliaprotocol.ai}).

\begin{itemize}
\tightlist
\item
  \textbf{Composition predicate and admission rule:}
  \texttt{packages/verify/quorum.js} (\texttt{verifyQuorum}) and
  \texttt{lib/signoff/quorum-session.js} (\texttt{canAccept},
  \texttt{quorumGate}).
\item
  \textbf{Conformance vectors:}
  \texttt{conformance/vectors/quorum.v1.json}, the eleven EP-QUORUM-v1
  vectors named in Section 4, each carrying real ES256 WebAuthn
  assertions. The JavaScript, Python, and Go reference verifiers
  (\texttt{packages/verify}, \texttt{packages/python-verify},
  \texttt{packages/go-verify}) are required to agree on every vector.
  These are cross-language reference verifiers, one team's ports of the
  same logic: a consistency check, not clean-room independent
  implementations. Independent implementations remain future
  interoperability evidence.
\item
  \textbf{Formal models:} \texttt{formal/\allowbreak ep\_handshake.tla}
  with \texttt{formal/\allowbreak ep\_handshake.cfg},
  \texttt{formal/\allowbreak ep\_quorum.als}, and
  \texttt{formal/\allowbreak ep\_relations.als}, with status and scope
  recorded in \texttt{formal/\allowbreak PROOF\_STATUS.md}.
\item
  \textbf{Runnable composition demo:} \texttt{node}
  \texttt{examples/multi-\allowbreak handshake/compose-\allowbreak and-verify.mjs}
  composes a 2-of-3 trail (threshold mode on the wire, with the declared
  order enforced by the composer at admission) that the reference
  verifier accepts, then demonstrates five refusals: an initiator
  self-approval, an out-of-order signature, a replayed member handshake
  (a stale or reused challenge), an off-roster key, and an off-roster
  member forced past admission that the verifier itself refuses (the
  executor does not trust the composer).
\item
  \textbf{Local end-to-end test:}
  \texttt{e2e/multi-party-quorum.spec.js} drives the full multi-party
  flow against a local development stack (it is opt-in and writes test
  data; it is a local test, not a production deployment). Three browser
  contexts with three CDP virtual authenticators enroll three approvers,
  an ordered program-officer, authorizing-official, inspector-general
  policy gates an initiator's action, consumption is refused with a 403
  \texttt{quorum\_not\_satisfied} after the first and second approvals,
  succeeds only after the third, and an out-of-order signer is refused
  before any approval is recorded.
\end{itemize}

\subsection{10. Conclusion}\label{conclusion}

The question ``M of N of what?'' is not an implementation detail; it
decides what the surviving artifact can prove. Composing signatures
answers ``enough keys participated,'' which is the right answer for
custody and the wrong one for accountability. Composing handshakes
(fresh, intent-bound, user-verified ceremonies with total outcomes)
answers the questions an irreversible action leaves behind: which
people, over exactly what, when, in what order, having demonstrably seen
it, and who refused. The composition preserves per-member freshness,
intent binding, attribution with refusal, order evidence, and separation
of duties precisely because it never collapses the members into one
artifact; its linear size is its evidentiary value. The safety core of
this process is machine-checked at the state-machine level and honestly
scoped, its predicate is pinned by adversarial conformance vectors, and
its limits (collusion, coercion, Sybil enrollment, semantic error) are
stated rather than waved away. For AI agents holding real operational
authority, we believe the unit of multi-party human control should be
the ceremony a human actually performed, and this paper is the argument
for building on that unit.

\subsection{References}\label{references}

{[}1{]} I. Schrock, ``Authorization Receipts for High-Risk Agent Actions
(EP),'' Internet-Draft draft-schrock-ep-authorization-receipts, 2026.

{[}2{]} I. Schrock, ``Multi-Party Quorum Authorization for High-Risk
Agent Actions (EP-QUORUM),'' Internet-Draft draft-schrock-ep-quorum,
2026.

{[}3{]} I. Schrock, ``The Two-Person Rule for AI Agents: Fail-Closed
Multi-Party Authorization with Offline-Verifiable Receipts,'' preprint,
Zenodo, DOI 10.5281/zenodo.20780638, 2026.

{[}4{]} C. Komlo and I. Goldberg, ``FROST: Flexible Round-Optimized
Schnorr Threshold Signatures,'' Selected Areas in Cryptography (SAC),
2020.

{[}5{]} ``The Flexible Round-Optimized Schnorr Threshold (FROST)
Protocol for Two-Round Schnorr Signatures,'' RFC 9591, 2024.

{[}6{]} R. Gennaro and S. Goldfeder, ``Fast Multiparty Threshold ECDSA
with Fast Trustless Setup,'' ACM CCS, 2018.

{[}7{]} J. Nick, T. Ruffing, and Y. Seurin, ``MuSig2: Simple Two-Round
Schnorr Multi-Signatures,'' CRYPTO, 2021.

{[}8{]} S. Micali, K. Ohta, and L. Reyzin, ``Accountable-Subgroup
Multisignatures,'' ACM CCS, 2001.

{[}9{]} R. Housley, ``Cryptographic Message Syntax (CMS),'' RFC 5652,
2009.

{[}10{]} A. Rundgren, B. Jordan, and S. Erdtman, ``JSON Canonicalization
Scheme (JCS),'' RFC 8785, 2020.

{[}11{]} W3C, ``Web Authentication: An API for accessing Public Key
Credentials, Level 2,'' W3C Recommendation, April 2021.

{[}12{]} U.S. Drug Enforcement Administration, ``Electronic
Prescriptions for Controlled Substances,'' 21 CFR Part 1311.

{[}13{]} V. Bertocci and B. Campbell, ``OAuth 2.0 Step Up Authentication
Challenge Protocol,'' RFC 9470, 2023.

\emph{Artifacts: the protocol drafts, the EP-QUORUM-v1 conformance
vectors, the formal models, and the reference verifiers are available in
the EMILIA Protocol open-source repository (Apache-2.0). Site:
\url{https://www.emiliaprotocol.ai}}

\end{document}
